\originalchapter{数据与抽样}
\label{ch:data}
\originalsection{数据、总体与统计问题}
防洪设施的高度、保险的保费和药物的疗效都需要处理尚未发生或无法全部观察的结果. 洪水高度会变化, 不同驾驶者的赔付也会变化; 在一组病人中看到疗效, 还不能保证下一组结果相同. MIT 原讲义用这些例子引入统计学. 数学模型把问题分成两部分: 可解释的结构与未解释的随机变化. 随机性既可能来自实际的随机机制, 也可能是复杂过程的近似描述.

\begin{definition}
总体是研究所针对的全部单位或结果, 样本是实际观察的单位或结果. 每一行数据对应一个观察单位, 每一列变量说明这个单位的一个特征. 参数是总体或模型的固定未知量; 统计量是仅由样本计算、不包含未知参数的函数.
\end{definition}
例如研究一个地区居民的收缩压, 总体可以是该地区在指定日期年龄在 18 岁以上的全部居民. 样本中的一行可以记录一位居民的年龄、性别和收缩压. 必须明确时间、对象和测量条件, 否则“平均血压”没有唯一对象. 以 $X_1,\ldots,X_n$ 表示样本测量值时, $\bar X=n^{-1}\sum_iX_i$ 是统计量, 总体均值 $\mu$ 是参数. 样本得到以前 $\bar X$ 是随机变量, 得到以后 $\bar x$ 是一个已知数值; 两者不能在置信陈述中混为一谈.

分类变量的数值编码并不意味着可以作所有算术. 给三种交通方式标为 1、2、3, 编码的均值不表达交通方式的平均程度. 有序变量允许比较先后, 但相邻等级的差值未必相等. 数值变量则需要区分离散计数与连续测量, 并保留单位. 数据表中的缺失项也不是零: 未回答收入与收入为零代表不同事实.

\begin{example}
两个独立公平骰子的点数和记为 $S$. 若猜中一个和可获 100 元, 应选择什么数? 若只知道一批药物试验中 100 位病人有 78 位好转, 要推断未来病人的好转概率, 又缺少什么?
\end{example}
\begin{solution}
骰子的机制已给定. 36 个有序点数组等可能, 和为 $s$ 的组合数为 $6-|s-7|$, 其中 $2\le s\le12$. 因而 $\PP(S=s)=(6-|s-7|)/36$, 选择 7 的获胜概率最大, 为 $1/6$. 这由已知分布算出结果概率, 属于概率问题.

药物问题的分布参数未知. 可以先用独立同分布的 Bernoulli 变量建模好转与否, 以样本比例 $0.78$ 估计好转概率 $p$. 但这个模型仍要求说明病人如何招募、何为好转、是否独立, 并需要量化估计误差. 如果要判断药物的因果效果, 还需要合适对照, 不能把有好转直接解释成药物有效. 后面的估计与检验处理在已声明模型条件下的量化问题.
\end{solution}

原课另用单个公平骰子的奖金说明期望与典型结果的区别: Alice 在点数至多为 3 时得 1 元, Bob 在点数至多为 2 时得 2 元. 期望收益分别为 $1/2$ 元与 $2/3$ 元, 若以期望收益为目标应选择 Bob. Alice 的获奖概率更高, Bob 的收益更分散. 期望收益、获奖频率与波动程度是不同指标.

\originalsection{抽样设计与偏差}
本节与后面的描述统计是自编补充. 参数推断中常假设独立同分布样本, 这一条件需要由数据取得方式支持. 对观察单位的选择发生系统性倾斜时, 增加样本量只会使错误目标估得更精确.

\begin{definition}
有限总体有 $N$ 个单位时, 简单随机不放回抽样指每个含 $n$ 个单位的子集被选择的概率都为 $\binom Nn^{-1}$. 分层抽样先把总体划为互不相交的层, 再在各层内分别抽样. 整群抽样随机选择若干群, 对选中的群进行调查或继续抽取单位.
\end{definition}
简单随机抽样中的“每个子集”比“每个单位入样概率相同”强. 后者并不能排除总是把某些单位绑在一起的方案. 不放回样本还不是严格的独立样本: 已抽到某人会影响下一个人的选择. 总体很大、抽样比例很小时可以讨论独立近似, 但不能把近似写成精确等式.

\begin{proposition}
有限总体数值为 $x_1,\ldots,x_N$, 总体均值 $\mu_N=N^{-1}\sum_jx_j$, 总体离差平方均值 $\sigma_N^2=N^{-1}\sum_j(x_j-\mu_N)^2$. 从中简单随机不放回抽取 $n$ 个值, 其均值 $\bar X$ 满足
\[
 \E\bar X=\mu_N,\qquad
 \Var(\bar X)=\frac{\sigma_N^2}{n}\frac{N-n}{N-1}\quad(N>1).
\]
\end{proposition}
\begin{proof}
按随机排列后的前 $n$ 个位置表示样本. 每个 $X_i$ 都均匀取总体中的一个值, 故均值为 $\mu_N$, 方差为 $\sigma_N^2$. 对 $i\ne k$, 有序的不相等单位等可能, 因而
\[
 \Cov(X_i,X_k)=\frac{1}{N(N-1)}\sum_{j\ne\ell}(x_j-\mu_N)(x_\ell-\mu_N)
 =-\frac{\sigma_N^2}{N-1}.
\]
最后一个等式用到全部离差之和为零: 不同项的乘积之和等于全部乘积和减去对角平方和. 将 $n$ 个方差与 $n(n-1)$ 个协方差相加, 再除以 $n^2$, 就得到所述公式.
\end{proof}
有限总体修正因子 $ (N-n)/(N-1)$ 在调查比例较大时显著降低方差; 当 $n=N$ 时全部单位已观察, 抽样误差为零. 若用通常无偏总体方差 $S_N^2=(N-1)^{-1}\sum_j(x_j-\mu_N)^2$, 同一式可写为 $(1-n/N)S_N^2/n$.

\begin{example}
某校有 900 位低年级与 100 位高年级学生, 各层平均每日自习时间分别为 2 小时和 4 小时. 调查时两层各选 50 人. 两层样本均值为 2.1 与 3.8 小时, 应如何估计全校均值?
\end{example}
\begin{solution}
总体目标是 $0.9\mu_1+0.1\mu_2$, 故分层估计为 $0.9(2.1)+0.1(3.8)=2.27$ 小时. 不加权平均 $(2.1+3.8)/2=2.95$ 小时将两层视为同样大, 估计了另一目标. 层内抽样均值各自无偏时, 按已知层比例加权的总体均值也无偏; 层内抽样若有系统性遗漏, 正确权重并不能消除它.
\end{solution}
自愿答卷、只调查容易接触的人、电话调查漏掉特定人群, 都可能形成选择偏差. 未回答与研究变量有关时会出现非应答偏差. 问题措辞、记忆错误和仪器误差属于测量问题. 这些问题需要通过设计、追访、校准及明确报告处理, 不是把标准误差写得更小就能解决.

抽样与随机分配解决不同问题. 随机抽样支持从样本推广到总体; 随机分配处理与对照帮助比较因果效果. 观察性研究中, 年龄同时影响治疗选择与疗效时, 两组疗效的差异可能混有年龄作用. 回归调整可以在模型条件下比较同类对象, 却不自动排除未观测混杂. 同一病人的多次测量也不能当成多个独立病人.

\originalsection{数据的中心与离散程度}
给定实数数据 $x_1,\ldots,x_n$, 排序后记为 $x_{(1)}\le\cdots\le x_{(n)}$. 均值使用每个值, 中位数按位置刻画中心. 两者回答的不是完全相同的问题.

\begin{definition}
样本均值为 $\bar x=n^{-1}\sum_ix_i$. 奇数 $n=2m+1$ 时中位数取 $x_{(m+1)}$, 偶数 $n=2m$ 时取 $(x_{(m)}+x_{(m+1)})/2$. 当 $n\ge2$, 样本方差与样本标准差分别为
\[
 s^2=\frac{1}{n-1}\sum_{i=1}^n(x_i-\bar x)^2,\qquad s=\sqrt{s^2}.
\]
\end{definition}
方差使用平方单位, 标准差与原数据同单位. 对 $y_i=a+bx_i$, 有 $\bar y=a+b\bar x$ 与 $s_y=|b|s_x$. 所以更换计量单位会改变数值, 不能把不同单位的方差直接相加比较.

\begin{proposition}
均值唯一最小化平方离差和 $\sum_i(x_i-a)^2$. 所有中位数最小化绝对离差和 $\sum_i|x_i-a|$; 偶数样本时区间 $[x_{(m)},x_{(m+1)}]$ 内的值全部取到最小值.
\end{proposition}
\begin{proof}
由于 $\sum_i(x_i-\bar x)=0$, 平方和展开为
\[
 \sum_i(x_i-a)^2=\sum_i(x_i-\bar x)^2+n(a-\bar x)^2,
\]
从而只有 $a=\bar x$ 最小. 对绝对离差, 将排序数据两端配成对. 每一对 $u\le v$ 满足 $|u-a|+|v-a|\ge v-u$, 当且仅当 $a\in[u,v]$ 等号成立. 偶数情形这些区间的交为中间两值的闭区间; 奇数情形还有中间项 $|x_{(m+1)}-a|$, 必须在中间值为零才能最小.
\end{proof}
这种优化性质说明, 单个极端值会对均值的平方误差目标产生很大影响, 而中位数主要由排序位置决定. “稳健”并不意味着中位数对所有统计目标都比均值好; 若目标就是总体期望, 仍要考虑均值及其估计误差.

\begin{proposition}
若 $X_1,\ldots,X_n$ 独立同分布, $\E X_i=\mu$ 且 $\Var(X_i)=\sigma^2<\infty$, 则 $\E\bar X=\mu$, $\Var(\bar X)=\sigma^2/n$, 而 $\E S^2=\sigma^2$.
\end{proposition}
\begin{proof}
均值与方差结论分别用期望线性性与独立性. 对样本方差, 恒等式
\[
 \sum_i(X_i-\bar X)^2=\sum_i(X_i-\mu)^2-n(\bar X-\mu)^2
\]
给出期望 $n\sigma^2-n(\sigma^2/n)=(n-1)\sigma^2$, 除以 $n-1$ 即可. 当不独立或总体方差不存在时, 此证明中均值方差的独立性计算或有限期望条件失效, 不能照用结论.
\end{proof}
标准差 $S$ 衡量单个数据的离散程度, 均值的估计标准误差 $S/\sqrt n$ 衡量均值在重复独立抽样中的变化. 调查报告中的“平均值加减一个标准差”与“平均值的置信区间”因此不是同一个对象.

\originalsection{分位数、直方图与箱线图}
\begin{definition}
经验分布函数为 $F_n(t)=n^{-1}\sum_i\ind_{\{x_i\le t\}}$. 对 $0<u\le1$, 经验分位数的一种约定为 $Q_n(u)=\inf\{t:F_n(t)\ge u\}=x_{(\lceil nu\rceil)}$. 第一、第三四分位数分别记为 $Q_1,Q_3$, 四分位距为 $\mathrm{IQR}=Q_3-Q_1$.
\end{definition}
分位数插值方式没有唯一的统一约定. 本节画箱线图时采用“先取中位数, 再对上下半组取中位数”的常用教学约定; 奇数样本的中间值不进入两半组. 它与上面的经验逆分位数定义不一定相同. 数据量小时差异尤其可见, 软件输出须报告所用约定.

直方图先将数轴分为若干区间. 等宽时柱高可以用频数或频率; 不等宽时若要用面积表达概率, 应取柱高为“该组相对频率除以组距”. 每根柱子的面积才等于该组相对频率, 全图面积为 1. 只比较不等宽柱子的高度可能产生错误印象. 类别变量则用条形图; 条形次序与间隔没有连续数轴的含义.

箱线图的箱体从 $Q_1$ 到 $Q_3$, 中间线为中位数. 一种常用须线规则是: 下须延伸至不低于 $Q_1-1.5\mathrm{IQR}$ 的最小观测值, 上须延伸至不高于 $Q_3+1.5\mathrm{IQR}$ 的最大观测值, 超出须线的观测单独标记. 须端不一定等于理论围栏, 也不一定是数据的最小值与最大值. 超出围栏只是在此规则下的异常点, 不证明数据错误或来自另一总体.

\begin{example}
十个观测值为 $1,2,3,4,5,6,7,8,9,20$. 计算均值、中位数、四分位距、样本标准差与箱线图须线, 并说明把最大值 20 换为 200 的影响.
\end{example}
\begin{solution}
总和为 65, 故均值为 6.5, 中位数为 5.5. 上下半组的中位数分别为 3 与 8, 所以 $\mathrm{IQR}=5$. 围栏为 $-4.5$ 和 $15.5$, 须端为 1 与 9, 20 单独标记. 平方和为 $1^2+\cdots+9^2+20^2=285+400=685$, 因而
\[
 s^2=\frac{685-10(6.5)^2}{9}=\frac{262.5}{9}=\frac{175}{6},\qquad s\approx5.401.
\]
若最大值换为 200, 均值变为 24.5, 而中位数和两四分位数保持原值. 新方差为 $(40285-10(24.5)^2)/9=34282.5/9$, 标准差约为 61.718. 这里中心的均值、中位数与离散指标呈现明显不同的敏感性.
\end{solution}
\begin{center}
\begin{tikzpicture}[x=.32cm,y=.65cm]
\draw[-{Stealth}] (0,0)--(22,0);
\foreach \x in {0,5,10,15,20} {\draw (\x,-.1)--(\x,.1);\node[below] at (\x,-.1){$\x$};}
\draw (1,.8)--(3,.8) (8,.8)--(9,.8);
\draw (1,.5)--(1,1.1) (9,.5)--(9,1.1);
\draw (3,.4) rectangle (8,1.2);
\draw (5.5,.4)--(5.5,1.2);
\fill (20,.8) circle (2pt);
\node[above] at (5.5,1.3){中位数 $5.5$};
\node[above] at (20,1){$20$};
\end{tikzpicture}
\end{center}
这是本书自编数据的 TikZ 箱线图. 同一数据如果使用经验逆分位数, $Q_n(.25)=x_{(3)}=3$, $Q_n(.75)=x_{(8)}=8$, 恰与半组约定相同; 换一组样本量便不必如此.

分布的偏斜、双峰与极端尾部常常不能由一个均值和方差完整表达. 画经验分布函数、直方图或散点图可以发现这些结构, 再决定模型是否合理. 图形也需要报告样本量、单位、分组方式和缺失数据处理.

\originalsection{相关、配对与初步推断}
对配对数据 $(x_i,y_i)$, 定义 $S_{xx}=\sum_i(x_i-\bar x)^2$, $S_{yy}=\sum_i(y_i-\bar y)^2$, $S_{xy}=\sum_i(x_i-\bar x)(y_i-\bar y)$. 当 $S_{xx},S_{yy}>0$ 时, 样本相关系数为 $r=S_{xy}/\sqrt{S_{xx}S_{yy}}$.

\begin{proposition}
样本相关系数满足 $-1\le r\le1$. 等号 $|r|=1$ 当且仅当两个非零中心化向量成比例, 即所有数据落在斜率非零的一条直线上.
\end{proposition}
\begin{proof}
对向量 $(x_i-\bar x)_{i=1}^n$ 与 $(y_i-\bar y)_{i=1}^n$ 使用 Cauchy--Schwarz 不等式得 $S_{xy}^2\le S_{xx}S_{yy}$. 该不等式取等当且仅当两个向量线性相关. 两者均非零时比例系数非零, 平移回去即为一条直线.
\end{proof}
$r$ 只度量线性关联. 对关于零对称的数据若 $y=x^2$, 可能有 $r=0$, 但 $y$ 完全由 $x$ 决定. 聚合不同群体还可能改变相关方向, 因而画散点图并检查分组比只报 $r$ 更有信息. 若每位病人都有治疗前后测量, 分析差值 $D_i=Y_{i,\mathrm{after}}-Y_{i,\mathrm{before}}$ 可以保留配对结构; 把两列作为独立样本会抹去同一病人的关联.

已知每位病人独立地以概率 $0.8$ 好转时, $K\sim\mathrm{Bin}(100,0.8)$, 所以 $\E K=80$, 而
\[
 \PP(K\ge65)=\sum_{k=65}^{100}\binom{100}{k}0.8^k0.2^{100-k}
 \approx0.99985286.
\]
原课把它写作约 $99.99\%$, 精确求和对应约 $99.9853\%$. 这是给定模型参数后预测样本结果的概率陈述. 换成观测结果反推未知参数, 计算方向便改变了.

MIT 原课中的 100 人有 78 人好转的例子, 用大样本正态近似构造
\[
 0.78\ \pm\ 1.96\sqrt{\frac{0.78(1-0.78)}{100}}
 \approx[0.6988,0.8612].
\]
此式估计的是模型中好转概率, 其“95\%”来自重复抽样时区间覆盖固定真值的比例, 并非已观测区间中每个数字有 95\% 的概率. 近似需要样本与概率不要过于接近边界, 也依赖独立同分布条件. 后文将推导它, 并给出不依赖正态近似的保守方法.

描述统计先忠实表达所见数据. 推断则需要由数据取得方式、概率模型和误差控制把结论延伸到未观察对象. 同一个样本比例既可写在描述表中, 也可进入参数估计与检验; 区别在于后一种用途必须说明概率保证来自何种模型.
